\documentclass{article}

\usepackage{corl_2026} % Use this for the initial submission.
% \usepackage[final]{corl_2026} % Uncomment for the camera-ready ``final'' version.
%\usepackage[preprint]{corl_2026} % Uncomment for pre-prints (e.g., arxiv); This is like ``final'', but will remove the CORL footnote.
% \documentclass{corl_2026}
\usepackage{graphicx}
\usepackage{amsmath}
\usepackage{amssymb}
\usepackage{booktabs, multirow}
\usepackage{subcaption}
\usepackage{pifont}
\usepackage[table]{xcolor}
\usepackage{pgfplots}
\pgfplotsset{compat=1.18}
\usepackage{algorithm}
\usepackage{algorithmic}
\usepackage[table]{xcolor}
\title{WAM-Nav: Asymmetric Latent World-Action Modeling for Unified Visual Navigation}

% The \author macro works with any number of authors. There are two
% commands used to separate the names and addresses of multiple
% authors: \And and \AND.
%
% Using \And between authors leaves it to LaTeX to determine where to
% break the lines. Using \AND forces a line break at that point. So,
% if LaTeX puts 3 of 4 authors names on the first line, and the last
% on the second line, try using \AND instead of \And before the third
% author name.

% NOTE: authors will be visible only in the camera-ready and preprint versions (i.e., when using the option 'final' or 'preprint'). 
% 	For the initial submission the authors will be anonymized.

\author{
  Ning Yang\\
  School of Intelligent Science and Technologys\\
  Nanjing University
  China\\
  \texttt{602025710040@smail.nju.edu.cn} \\
  %% examples of more authors
  %% \And
  %% Coauthor \\
  %% Affiliation \\
  %% Address \\
  %% \texttt{email} \\
  %% \AND
  %% Coauthor \\
  %% Affiliation \\
  %% Address \\
  %% \texttt{email} \\
  %% \And
  %% Coauthor \\
  %% Affiliation \\
  %% Address \\
  %% \texttt{email} \\
  %% \And
  %% Coauthor \\
  %% Affiliation \\
  %% Address \\
  %% \texttt{email} \\
}


\begin{document}
\maketitle
%===============================================================================
\begin{abstract}
    % Visual navigation requires generating smooth and collision-free trajectories under complex geometric and physical constraints. Existing reactive policies that directly map partial observations to actions often fail in long-horizon scenarios due to the lack of predictive reasoning. While visual world models provide predictive foresight, 但策略与想象相互独立训练的conventional 模块化解耦 pipelines suffer from error accumulation between future prediction and action execution破坏了导航轨迹的一致性.为此，
    % We propose \textbf{WAM-Nav}, a first \textbf{Latent World-Action Model} for embodied visual navigation 通过联合训练策略和想象 jointly models action generation and future visual prediction，在不降低推理效率的同时提升了导航决策性能. WAM-Nav employs a shared Diffusion Transformer in a unified latent space to asymmetric joint diffusion generate action trajectories and latent visual consequences，避免了多步自回归预测带来的推理延迟和视觉误差累积.
    % To improve long-horizon 轨迹 consistency, we introduce a dual-stream contextual condition in WAM-Nav that incorporates an episode-level ego-motion history with sequential visual observations. Corporating with a multi-modal goal alignment mechanism, the WAM-Nav is able to seamlessly support image-goal, point-goal, and no-goal navigation.
    % Experiments on the challenging \textit{ClutterScenes} and \textit{InternScenes} benchmarks show that WAM-Nav achieves superior generalization,特别是在Image-goal和point-goal导航任务上分别提升15.7\%，3.3\%. Real-world deployment further demonstrates effective 以85\%的任务成功率进一步证实了其zero-shot sim-to-real transfer.
    Visual navigation requires generating smooth and collision-free trajectories under complex geometric and physical constraints. Existing reactive policies that directly map observations to actions lack anticipatory reasoning, limiting their ability to proactively avoid obstacles. While visual imagination offers predictive foresight, conventional modular approaches separate scene prediction from policy learning, often leading to error accumulation and inefficient inference.
    To address these limitations, we propose WAM-Nav, a Latent World-Action Model for embodied visual navigation that jointly learns action generation and latent future visual prediction, enabling more robust and foresighted navigation decisions without compromising inference efficiency.
    Specifically, WAM-Nav utilizes a shared Diffusion Transformer for asymmetric joint diffusion to concurrently generate long-horizon actions and short-horizon visual foresight, reducing the inference latency and visual error accumulation inherent in multi-step autoregressive rollouts.
    To further encourage smooth and consistent trajectory generation, we introduce a dual-stream contextual conditioning mechanism that integrates episode-level ego-motion history with sequential visual observations. Combined with a unified goal alignment module that preserves balanced representations across goal types, WAM-Nav naturally supports image-goal, point-goal, and no-goal navigation within a single policy.
    Extensive experiments on the challenging \textit{ClutterScenes} and \textit{InternScenes} benchmarks demonstrate strong generalization of WAM-Nav, particularly on image-goal and point-goal navigation, where it improves success rates by \textbf{15.7\%} and \textbf{3.3\%}, respectively. Real-world deployment further validates effective zero-shot sim-to-real transfer, achieving an average \textbf{85\%} task success rate across diverse indoor and outdoor environments.
\end{abstract}

% Two or three meaningful keywords should be added here
\keywords{Embodied Visual Navigation, World Action Model, Diffusion Transformer} 

% -----------------------------------------------------------------------------
% WAM-Nav Introduction Section (Concise Version for 2-Page Fit)
% -----------------------------------------------------------------------------

\section{Introduction}

Embodied visual navigation \citep{wei2026navol, yang2025transformer, yang2025longterm, zhang2025navfom} aims to guide agents through complex, unseen physical environments by generating smooth, collision-free trajectories based on visual observations \citep{zhu2021survey}. Over the past few years, the dominant navigation paradigm has gradually transitioned from traditional decoupled mapping-and-planning pipelines \citep{campos2021orb, labbe2019rtabmap, chaplot2020learning} to learning-based intelligent methods. As shown in Figure \ref{fig:intro} (a), recent end-to-end reactive approaches \cite{shah2023gnm, shah2023vint, sridhar2024nomad, peng2025logoplanner, cai2025navdp} directly map visual observations to executable actions. 
Moreover, to endow agents with predictive reasoning capabilities, recent research has actively explored visual imagination and world-model-based navigation strategies. These modular approaches generally fall into two distinct paradigms. The first paradigm prioritizes generating future visual subgoals, which are subsequently utilized to infer control actions via inverse dynamics models \cite{qin2025navigatediff, ni2024generatesubgoal}. The second paradigm focuses on trajectory evaluation, where candidate paths are first generated and then scored by a world model that imagines future visual states\cite{bar2025nwm,zhang2026raenwm, dong2025uniwm}. 

Although existing approaches have improved robotic navigation to varying degrees, fundamental limitations remain across current paradigms. Reactive end-to-end mapping methods, which directly map observations to actions, rely heavily on instantaneous perception and therefore lack predictive reasoning and local dynamics understanding, making them prone to local optima and collisions in cluttered environments. Modular decoupled approaches provide explicit foresight, but their separate training of action decision-making and future imagination introduces substantial computational latency and compounding errors, often resulting in delayed responses and poor trajectory consistency. While existing world-action models have shown promising joint visual-action modeling capabilities in robot manipulation \citep{ye2026dreamzero, bi2025motus, antgroup2026lingbotva}, their autoregressive generation paradigm remains challenging for navigation due to limited real-time capability and accumulated prediction errors under large viewpoint changes. Furthermore, most existing navigation methods are designed for a single target specification \citep{shah2023gnm,shah2023vint,sridhar2024nomad,peng2025logoplanner,bar2025nwm,qin2025navigatediff,dong2025uniwm}, requiring redesign, new data collection, and retraining when adapting to new tasks. Although NavDP \citep{cai2025navdp} supports multiple goal-oriented navigation tasks, its unimodal alignment design still leads to imbalanced performance across task types.

% 尽管上述方法都在一定程度上提升了机器人的导航能力，但是reactive 端到端映射方法due to reliance on instantaneous perception, these reactive methods lack predictive reasoning and local dynamics understanding, making them prone to local optima and collisions in cluttered environments with complex obstacle configurations. 同时，While模块化解耦方法providing valuable foresight, 这种动作决策与未来想象分离独立训练的方法进一步引入了introduces substantial computational latency and compounding errors，造成导航卡顿和轨迹一致性差。现有世界动作模型在机器人操作中已展示出联合建模视觉与动作的潜力 \citep{ye2026dreamzero, bi2025motus, antgroup2026lingbotva}，但对导航任务而言这类自回归逐帧生成范式不仅难以满足导航任务的实时性要求，而且会由于导航视角变动大造成未来帧预测累积误差变大进而错误引导决策。另外，大多仅支持一种目标类型的导航方法 \citep{shah2023gnm,shah2023vint,sridhar2024nomad,peng2025logoplanner,bar2025nwm,qin2025navigatediff,dong2025uniwm}当更换导航任务时需要重新更改模型，收集数据retrain，即使NavDP \citep{cai2025navdp}可支持多种目标类型导航任务，但是由于单模态对齐，导致不同任务之间的模型能力表现不平衡。

\begin{figure}[t]
    \centering
    \includegraphics[width=\linewidth]{corl_2026_template_cameraready/figures/Introduction.png}
    \caption{Overview of the \texttt{WAM-Nav} paradigm and performance.(a) Methodological comparison: Unlike traditional purely reactive or decoupled modular pipelines, \texttt{WAM-Nav} jointly models action generation and latent future visual prediction within a unified framework. (b) Quantitative performance: Our method achieves competitive results against established baselines across diverse evaluation settings.}
    \label{fig:intro}
    \vspace{-5mm}
\end{figure}

Motivated by these advances, we propose \texttt{WAM-Nav}, a latent world-action model for embodied visual navigation that jointly learns action generation and latent future visual prediction, enabling more robust and foresighted navigation decisions without compromising inference efficiency, as illustrated in Figure \ref{fig:intro}. Instead of relying on sequential decoupled modules, \texttt{WAM-Nav} utilizes a shared Diffusion Transformer (DiT) for asymmetric joint diffusion to concurrently generate long-horizon actions and short-horizon visual foresight. This joint modeling reduces the inference latency and visual error accumulation inherent in multi-step autoregressive rollouts, while ensuring a deep coupling between visual features and control dynamics to empower the policy to anticipate environmental variations and proactively steer clear of obstacles. To further guide generation and preserve trajectory consistency, we introduce Dual-Stream Contextual Conditioning (DSCC), which integrates episode-level ego-motion history with sequential visual observations. Combined with a unified goal alignment module that preserves balanced representations across goal types, \texttt{WAM-Nav} naturally supports image-goal, point-goal, and no-goal navigation within a single policy without complex architecture switching.

We conduct extensive evaluations on cluttered and physically realistic benchmarks, including the \textit{ClutterScenes} and \textit{InternScenes}. Experimental results demonstrate that \texttt{WAM-Nav} achieves strong zero-shot generalization, particularly on image-goal and point-goal navigation where it improves success rates by \textbf{15.7\%} and \textbf{3.3\%}, respectively. Furthermore, deployment on physical robot platforms validates its capability for sim-to-real transfer, achieving an average \textbf{85\%} task success rate across diverse indoor and outdoor environments.


% -----------------------------------------------------------------------------
% Contributions Section
% -----------------------------------------------------------------------------
% In summary, our main contributions are as follows:
% \begin{itemize}
%     \item We propose \texttt{WAM-Nav}, an end-to-end framework that deeply couples action prediction and visual imagination within a compact latent space via a shared DiT, mitigating the latency and compounding errors of decoupled systems.
%     \item We introduce the Dual-Stream Prior Representation (DSCC), fusing historical ego-motion with sequential visuals to improve the physical consistency of generated trajectories.
%     \item We design a multi-modal goal projection mechanism enabling a single policy to flexibly support heterogeneous navigation tasks, validated by extensive simulated and real-world experiments.
% \end{itemize}

\section{Related Work}

\subsection{Learning-based Embodied Visual Navigation}

In contrast to traditional mapping-and-planning methods, the end-to-end learning framework has emerged as a dominant paradigm for embodied visual navigation \citep{zhu2021survey, zhu2017target, savva2019habitat, wijmans2020ddppo, zeng2025navidiffusor}. Recent foundation models have significantly advanced this field: GNM \citep{shah2023gnm} demonstrated cross-embodiment generalization using a lightweight CNN policy \citep{krizhevsky2012imagenet}, while ViNT \citep{shah2023vint} extended navigation horizons via Transformer-based sub-goal planning. Subsequently, NoMaD \citep{sridhar2024nomad} introduced diffusion policies \citep{janner2023policy} to model multi-modal action distributions, and NavDP \citep{cai2025navdp} incorporated an actor-critic mechanism to score and select the safest generated trajectories. Despite these advances, such reactive policies directly map local observations to actions without explicitly modeling their future perceptual consequences, rendering them susceptible to collisions in cluttered environments \citep{qin2025navigatediff, shen2026efficient}. In contrast, \texttt{WAM-Nav} addresses this limitation by integrating predictive visual foresight directly into the policy, enabling proactive obstacle avoidance and safer navigation decision-making.

\subsection{World Models and Visual Imagination for Navigation}

Integrating visual imagination facilitates foresight-driven navigation \citep{zhang2026sparsevideonav, nie2025wmnav}, and early systems typically adopt modular architectures to evaluate candidate motions. For instance, Pathdreamer \citep{koh2021pathdreamer} and NWM \citep{bar2025nwm} rely on decoupled visual prediction and trajectory planning, synthesizing future observations via a world model built atop an existing navigation policy. Similarly, NavigateDiff \citep{qin2025navigatediff} employ diffusion models \citep{ho2020denoising} to imagine future visual sub-goals, which are subsequently executed by a separate inverse dynamics policy. These approaches treat visual imagination and trajectory generation as sequential and decoupled processes. This separation frequently introduces state-action misalignment and compounding control errors. To address these misalignments, frameworks like UniWM \citep{dong2025uniwm} attempt to couple planning with foresight through a unified world model. However, their reliance on explicit pixel-level generation introduces substantial computational latency and dedicates model capacity to rendering irrelevant background details rather than optimizing core decision-making. Consequently, predicting future states in compact feature spaces via latent imagination \citep{hafner2020dreamer, clam2025, navmorph2025, beingh072025} offers a more efficient alternative for embodied control. To capitalize on this, \texttt{WAM-Nav} jointly models action generation and visual foresight, where future visual dynamics are predicted in a latent space. Rather than decoupling prediction from decision-making, it learns both jointly to promote temporal consistency and physical feasibility with minimal inference latency.

% \section{Problem Definition}
% \label{sec:problem}

% We consider the task of mapless visual navigation under three distinct goal modalities: image-goal ($g^{\mathrm{img}}$), point-goal ($g^{\mathrm{pt}}$), and goal-agnostic free exploration ($g^{\varnothing}$). At any given time step $t$ within a partially observable, unseen environment, the agent receives a visual history $\mathcal{O}_t=\{o_{t-k+1}, \ldots, o_t\}$, an action-accumulated ego-motion history $\mathcal{S}_t=\{s_{t-k+1}, \ldots, s_t\}$, and the specified goal condition $g \in \{g^{\mathrm{img}}, g^{\mathrm{pt}}, g^{\varnothing}\}$. 

% To transition from myopic reactive control to foresight-driven navigation, we formulate this task as a conditional joint generation problem. The policy simultaneously predicts a future action trajectory $\mathbf{A}_t=\{a_t,\ldots,a_{t+H-1}\}$ and its corresponding latent visual consequences $\mathbf{Z}_{t+1:t+H}=\{z_{t+1}, \ldots, z_{t+H}\}$. The learning objective is to model the joint distribution:
% \begin{equation}
%     p_\theta(\mathbf{A}_t, \mathbf{Z}_{t+1:t+H} \mid \mathcal{O}_t, \mathcal{S}_t, g).
% \end{equation}
% Through this joint formulation, the policy leverages the historical context $(\mathcal{O}_t, \mathcal{S}_t)$ to generate an action sequence $\mathbf{A}_t$ that remains consistent with the predicted future states $\mathbf{Z}_{t+1:t+H}$, ultimately navigating the agent safely toward the goal $g$.
%
\section{Methodology}
\label{sec:method}
% \subsection{Overview}
% 如图2所示，\texttt{WAM-Nav}整体框架中主要包括unified goal alignment,Dual-Stream Contextual Conditioning和Asymmetric Action-Foresight Generation三部分。首先，为平衡应对多种导航任务，我们将每一种目标类型进行统一双向对齐；其次，构建包含机器人历史RGB视觉观测和episode-level的ego-motion trajectory的双流上下文，并经分别显式注入目标信息之后进行压缩融合；最后，以此上下文为condition，在shared DiT中异步联合去噪生成动作轨迹和latent视觉想象。
As illustrated in Figure~\ref{fig:pipeline}, \texttt{WAM-Nav} consists of three key components: (1) \textit{Unified Goal Alignment}, (2) \textit{Dual-Stream Contextual Conditioning}, and (3) \textit{Asymmetric Action-Foresight Generation}. Specifically, to enable balanced performance across diverse navigation tasks, we first project heterogeneous goal specifications into a unified alignment space. We then construct a dual-stream spatiotemporal context by jointly encoding the robot’s sequential visual observations and episode-level ego-motion history. Both streams are explicitly modulated by goal-relevant information and aggregated into a compact conditioning representation $C$. Conditioned on $C$, a shared DiT jointly generates future action trajectories and latent visual foresight through asymmetric denoising.


\begin{figure}[t]
    \centering
    \includegraphics[width=\linewidth]{corl_2026_template_cameraready/figures/Method.png}
    \caption{Architecture overview of the \texttt{WAM-Nav} framework. Heterogeneous navigation goals are explicitly routed into visual-semantic ($g_V$) and trajectory geometric ($g_G$) queries. These queries contextually modulate historical RGB-D sequences and relative ego-motion trajectories, synthesizing a compact context condition $C$. Conditioned on $C$, a shared DiT performs asymmetric joint generation of future control actions and latent visual foresight.}
    \label{fig:pipeline}
    \vspace{-5mm}
\end{figure}

\subsection{Unified Goal Alignment}

To enable balanced performance across Image-Goal, Point-Goal, and No-Goal exploration within a single policy, \texttt{WAM-Nav} introduces a unified goal alignment mechanism. Unlike prior navigation approaches~\citep{cai2025navdp, wang2025internvla} that reformulate all tasks as point-goal navigation, our design preserves modality-specific goal information to maintain task-specific expressiveness across different navigation settings. Given a goal input $g$, \texttt{WAM-Nav} encodes it into two complementary goal embeddings: a visual semantic query $g_V \in \mathbb{R}^D$ for visual memory retrieval, and a geometric query $g_G \in \mathbb{R}^D$ for trajectory-level directional guidance.

Formally, the goal input $g$ is first transformed into a base embedding $e_g$ via a modality-specific feature extractor $E_{\phi}(\cdot)$, and then projected into the two functional token spaces:
\begin{equation}
g_V = \psi_V(e_g), \quad g_G = \psi_G(e_g), \quad \text{with} \quad e_g = E_{\phi}(g)
\end{equation}
where $E_{\phi}(\cdot)$ is instantiated according to the goal modality, including a Vision Transformer trained from scratch for image goals, sinusoidal positional encoding for relative coordinates, and a masked zero-state representation for goal-free exploration. The projection operators $\psi_V(\cdot)$ and $\psi_G(\cdot)$ are learnable linear mappings that map the base embedding into the visual semantic and geometric spaces, respectively.

This alignment design preserves modality-specific goal information while providing a unified interface across heterogeneous navigation tasks, enabling balanced performance across different goal settings and producing goal-conditioned context for the subsequent Section~\ref{sec:dscc}, where it serves as conditioning for the joint modeling of action and visual foresight.

\subsection{Dual-Stream Contextual Conditioning}
\label{sec:dscc}

While visual history provides spatial cues for obstacle avoidance, relying solely on visual conditioning often leads to kinematically inconsistent or jittery trajectories due to the lack of explicit momentum constraints. To address this, \texttt{WAM-Nav} introduces Dual-Stream Contextual Conditioning (DSCC) over a historical sliding window from $t-k+1$ to $t$. We dynamically construct the conditioning context $C$ by fusing a goal-modulated visual memory stream (capturing local spatial geometry) with a trajectory-aware kinematic history stream (capturing physical momentum). Fusing these streams provides the shared DiT with a spatiotemporal conditioning context that ensures the generated action trajectories are both geometrically collision-free and kinematically smooth.

\textbf{Goal-Modulated Visual Memory.} 
Historical RGB observations $\mathcal{O}_t$ are processed by DINO-v2 \citep{oquab2023dinov2} into a memory tensor $V \in \mathbb{R}^{B \times (k \cdot P) \times D}$, where $B$ represents the batch size, $k$ is the historical window length, and $P$ is the patch count per frame. To selectively reinforce spatial tokens relevant to the target, the visual query $g_V \in \mathbb{R}^{B \times D}$ computes independent relevance scores for each patch token via a scaled dot-product:
\begin{equation}
    \alpha = \sigma \left( \frac{g_V V^\top}{\sqrt{D}} \right) \in \mathbb{R}^{B \times (k \cdot P) \times 1}
\end{equation}
where $\sigma$ is the sigmoid function. The visual memory is then residually updated:
\begin{equation}
    \tilde{V} = V + \alpha \odot V
\end{equation}

\textbf{Trajectory-Aware Motion History.} 
To incorporate the agent's kinematic momentum and ensure robust zero-shot generalization, the motion trajectory stream processes the executed pose sequence $\mathcal{S}_t$. Absolute poses are converted into coordinate-invariant relative displacements and heading changes in the current egocentric frame, yielding a relative kinematics sequence $\tilde{\mathcal{S}}_t = \{(\Delta x_i, \Delta y_i, \Delta \theta_i)\}_{i=t-k+1}^{t}$ (detailed in Appendix~\ref{appendix:ego_history_conversion}). 

This relative motion sequence $\tilde{\mathcal{S}}_t$ is embedded and processed by a causal Transformer encoder $H \in \mathbb{R}^{B \times k \times D}$ to capture temporal kinematic profiles. To align past momentum with current target directionality, we query $H$ with the geometric goal embedding $g_G \in \mathbb{R}^{B \times D}$ via cross-attention to extract a consolidated kinematic vector $o_{\mathrm{kin}}$:
\begin{equation}
    o_{\mathrm{kin}} = \mathrm{CrossAttn}(g_G, H, H) \in \mathbb{R}^{B \times D}
\end{equation}
The vector $o_{\mathrm{kin}}$ mathematically summarizes the agent's historical motion continuity relative to its target directionality.

\textbf{Condition Fusion via Cross-Attention.} 
To structurally enforce kinematic momentum to guide spatial visual retrieval, the visual and trajectory streams are integrated via a multi-layer Transformer Decoder to formulate the unified conditioning context $C \in \mathbb{R}^{B \times N_c \times D}$. Specifically, the extracted kinematic token $o_{\mathrm{kin}}$ is projected to bias a set of learnable query embeddings $Q_c \in \mathbb{R}^{N_c \times D}$:
\begin{equation}
    C = \mathrm{TransformerDecoder} \left( Q_c + \phi(o_{\mathrm{kin}}), \tilde{V}, \tilde{V} \right)
\end{equation}
where $\phi(\cdot)$ denotes a linear projection. Architecturally, this fusion mechanism utilizes the agent's kinematic momentum ($o_{\mathrm{kin}}$) to actively query the goal-modulated visual spatial context ($\tilde{V}$). The resulting condition $C$ thus explicitly constrains the subsequent latent generation, enforcing both physical execution smoothness and geometric safety.

\subsection{Asymmetric Action-Foresight Generation}

Conditioned on the unified context condition $C$, \texttt{WAM-Nav} jointly models the future action trajectory $\mathbf{A}_t = \{a_t, \dots, a_{t+H_{\mathrm{act}}-1}\}$ over a long planning horizon $H_{\mathrm{act}}$ and its corresponding latent visual foresight $\mathbf{Z}_{t+1:t+H_{\mathrm{vis}}} = \{z_{t+1}, \dots, z_{t+H_{\mathrm{vis}}}\}$ over a short prediction horizon $H_{\mathrm{vis}}$ ($H_{\mathrm{vis}} \le H_{\mathrm{act}}$) within a shared representation space. The future visual state $z_i = \mathcal{E}(o_i)$ is compressed via a pre-trained Stable Diffusion VAE \citep{stabilityai2022sd} into a compact grid of $N$ latent patches. 

To optimize this generative process, we employ the flow-matching \citep{streamingflow2025} training strategy to learn the joint conditional distribution $p_\theta(\mathbf{A}, \mathbf{Z} \mid C)$ via asymmetric joint diffusion. This strategy formulates the continuous probability paths as straight lines that map standard Gaussian priors $(\mathbf{A}_0, \mathbf{Z}_0) \sim \mathcal{N}(\mathbf{0}, \mathbf{I})$ to the empirical data manifold. At any flow time $\tau \in [0, 1]$, the interpolated states are defined as:
\begin{equation}
    \mathbf{A}_{\tau} = (1-\tau)\mathbf{A}_0 + \tau \mathbf{A}_1, \quad
    \mathbf{Z}_{\tau} = (1-\tau)\mathbf{Z}_0 + \tau \mathbf{Z}_1
\end{equation}
where the corresponding target velocity fields are given by $u_A = \mathbf{A}_1 - \mathbf{A}_0$ and $u_Z = \mathbf{Z}_1 - \mathbf{Z}_0$.

The architectural backbone for this asymmetric generative progress is a shared DiT. The noised sequences $\mathbf{A}_{\tau}$ and $\mathbf{Z}_{\tau}$ are tokenized, concatenated, and processed through a multi-layer DiT. Crucially, within each DiT block, these heterogeneous tokens undergo a shared self-attention, enabling the dynamic exchange of spatiotemporal constraints between the physical action paths and the latent visual representations at every layer. Both streams cross-attend to the context condition $C$, with the timestep $\tau$ modulated via adaptive LayerNorm (adaLN) layers to regress the joint velocity fields:
\begin{equation}
    \hat{u}_A, \hat{u}_Z = f_{\theta}(\mathbf{A}_{\tau}, \mathbf{Z}_{\tau}, \tau, C).
\end{equation}
This shared parameterization makes latent foresight act as a perception-grounded constraint on action generation, penalizing action--scene inconsistency through the visual velocity-matching loss. Unlike manipulation-oriented WAMs~\citep{ye2026dreamzero, bi2025motus}, where future visual changes are often local and object-centric, navigation involves large egocentric viewpoint changes; long autoregressive visual rollouts would introduce both inference latency and accumulated visual errors that may misguide action generation. \texttt{WAM-Nav} therefore adopts an asymmetric design: long-horizon actions preserve trajectory continuity, while short-horizon latent foresight provides reliable near-future geometric constraints without overextending visual prediction.

\subsection{Model Training and Inference}

We train \texttt{WAM-Nav} end-to-end by minimizing a joint objective consisting of flow-matching velocity regression and multi-modal goal alignment:
\begin{equation}
    \label{eq:l_total}
    \mathcal{L}_{\mathrm{total}} = \mathbb{E}_{\tau, \mathbf{A}_0, \mathbf{Z}_0} \left[ \|\hat{u}_A - u_A\|_2^2 + \lambda_{\mathrm{img}}\|\hat{u}_Z - u_Z\|_2^2 \right] + \lambda_{\mathrm{align}}\mathcal{L}_{\mathrm{align}}
\end{equation}
where $\tau \sim \mathcal{U}(0,1)$, and $\lambda_{\mathrm{img}}, \lambda_{\mathrm{align}}$ are scalar weighting coefficients. The term $\mathcal{L}_{\mathrm{align}}$ represents a symmetric contrastive InfoNCE loss. By maximizing the mutual information between cross-space projections of equivalent objectives, this alignment constraint ensures consistency across diverse goal modalities, facilitating the extraction of unified directional and visual cues regardless of the original input format.

During online deployment, \texttt{WAM-Nav} operates within a receding-horizon control loop. At each execution step, conditioned on the current prior $C$, candidate trajectories are sampled by integrating the learned vector fields from a standard Gaussian prior using an Euler ODE solver over $N_{\mathrm{steps}}$ integration steps. Following standard operational protocols for generative navigation policies \citep{sridhar2023nomad}, the framework extracts the initial action segment from the generated trajectory to drive the robot's low-level executors. This receding-horizon procedure is continuously repeated, effectively balancing high-frequency closed-loop reactivity with long-horizon spatiotemporal foresight.

\section{Experimental Results}
\label{sec:result}
We conduct extensive zero-shot deployments of \texttt{WAM-Nav} to evaluate its effectiveness, focusing on five research questions: 
\textbf{Q1:} Does \texttt{WAM-Nav} achieve superior zero-shot generalization while maintaining balanced performance across diverse navigation tasks?  
\textbf{Q2:} What mechanisms enable \texttt{WAM-Nav} to navigate highly cluttered environments more effectively than existing baselines?  
\textbf{Q3:} How do the individual components contribute to the overall performance improvement?  
\noindent\textbf{Q4:} Can \texttt{WAM-Nav} transfer to physical platforms for zero-shot real-world navigation?
\subsection{Simulation Experiments}

\noindent\textbf{Training Datasets and Setup.}
Consistent with prior works \citep{peng2025logoplanner,cai2025navdp}, \texttt{WAM-Nav} is trained on the large-scale visual navigation dataset VLN-N1 \citep{wang2025internvla, interndata2025}. VLN-N1 is built upon six categories of 3D scene assets, including Replica~\citep{Straub2019Replica}, Matterport3D~\citep{Chang2017Matterport3D}, Gibson~\citep{Xia2018Gibson}, 3D-FRONT~\citep{fu20203dfront}, HSSD~\citep{Khanna2023HSSD}, and HM3D~\citep{Ramakrishnan2021HM3D}. Leveraging an extensive domain randomization pipeline, the dataset provides over 400 hours of collision-free and smooth first-person navigation trajectories collected in simulation, comprising more than 200K trajectories in total.

\texttt{WAM-Nav} is trained with a learning rate of \(1.5\times10^{-4}\). We use a per-GPU batch size of 256, resulting in an effective batch size of 2048 across 8 NVIDIA A800 GPUs. The total training cost is approximately \(8\times48\) GPU hours. Additional implementation details are provided in Appendix~\ref{appendix:implementation_details}.

\noindent\textbf{Zero-Shot Evaluation Platform.}
We evaluate \texttt{WAM-Nav} and all baselines in a zero-shot setting on a simulation benchmark built upon IsaacSim, using the wheeled robot ClearPath Dingo as the navigator. The benchmark includes two categories of test environments: \textit{Cluttered} scenes and \textit{InternScenes} \citep{cai2025navdp, wang2025internvla, wang2024grutopia}.
The \textit{Cluttered} scenes contain 10 \textit{easy} scenes and 10 \textit{hard} scenes. Both are constructed from randomly generated layouts with cluttered obstacles, while the \textit{hard} scenes feature denser obstacle distributions. The \textit{InternScenes} include 20 \textit{home} scenes and 20 \textit{commercial} scenes. \textit{Home} scenes are characterized by narrow passages and cluttered semantic layouts, whereas \textit{commercial} scenes cover representative indoor categories such as \textit{hospitals}, \textit{supermarkets}, and \textit{restaurants}. For each scene, we randomly sample 100 navigation episodes, yielding a total of 6,000 evaluation episodes. In every episode, the robot is initialized at a random location within the scene. The image-goal, point-goal, and no-goal settings are evaluated in the same environments, differing only in task objectives.

\noindent\textbf{Metrics.}
For image-goal and point-goal navigation, we evaluate \textbf{Success Rate (SR)} and \textbf{Success weighted by Path Length (SPL)} to measure task completion and navigation efficiency. For no-goal exploration, we report episode \textbf{Time} and explored \textbf{Area} to assess exploration efficiency, collision avoidance, and scene coverage. Higher values indicate better performance for all metrics.

\noindent\textbf{Compared Methods.}

To comprehensively evaluate the effectiveness of \texttt{WAM-Nav}, we benchmark our model against a diverse set of strong baselines tailored to different navigation tasks (detailed configurations are provided in Appendix~\ref{appendix:baseline_details}). Specifically, for \textbf{image-goal} navigation, we compare \texttt{WAM-Nav} against GNM~\citep{shah2023gnm}, ViNT~\citep{shah2023vint}, NoMaD~\citep{sridhar2024nomad}, NWM~\citep{bar2025nwm}, and NavDP~\citep{cai2025navdp}. For \textbf{point-goal} navigation, the baselines include DD-PPO~\citep{wijmans2020ddppo}, iPlanner~\citep{yang2023iplanner}, ViPlanner~\citep{roth2024viplanner}, and NavDP~\citep{cai2025navdp}. Finally, for \textbf{no-goal} exploration, we compare against GNM~\citep{shah2023gnm}, ViNT~\citep{shah2023vint}, NoMaD~\citep{sridhar2024nomad}, and NavDP~\citep{cai2025navdp}.

\noindent\textbf{Result Analysis.}
For \textbf{Q1}, Table~1 summarizes the zero-shot evaluation results across Image-Goal, Point-Goal, and No-Goal navigation. \texttt{WAM-Nav} achieves the best average performance across all three tasks, reaching 50.2\% SR / 48.2\% SPL on Image-Goal, 80.4\% SR / 78.0\% SPL on Point-Goal, and 168.9 m$^2$ explored area in No-Goal exploration. Notably, \texttt{WAM-Nav} consistently outperforms prior methods across both \textit{ClutterScenes} and \textit{InternScenes}, demonstrating strong out-of-domain generalization under varying scene layouts and semantic complexity. These results indicate that \texttt{WAM-Nav} provides a well-balanced general-purpose navigation policy without task-specific architectural adaptation.

For \textbf{Q2}, Figure~\ref{fig:result} qualitatively compares NavDP and \texttt{WAM-Nav} on the Image-Goal task. NavDP, as a reactive policy, often delays obstacle avoidance until close proximity, leading to abrupt trajectory shifts and unstable navigation. In contrast, \texttt{WAM-Nav} jointly models future trajectories and visual outcomes, enabling proactive obstacle avoidance and smoother path planning. By combining visual observations with historical motion context, \texttt{WAM-Nav} generates more consistent trajectories. Its decoded visual predictions also align closely with ground-truth observations, validating the effectiveness of the joint imagination mechanism in providing reliable spatial guidance for robust navigation.

For \textbf{Q3}, Table~\ref{tab:ablation_study_flattened_onedecimal_final} evaluates the roles of contextual motion trajectories and latent visual prediction. Latent prediction alone improves SR from 42.1\% to 45.7\%, with notable gains in \textit{ClutterScenes} (+6.4\%), highlighting its effectiveness for obstacle avoidance. In contrast, using only motion trajectories degrades performance in complex scenes (e.g., \textit{InternScenes}: 33.2\% $\rightarrow$ 30.0\%), suggesting that historical motion alone lacks sufficient foresight for timely adaptation. Combining both achieves the best performance (50.2\% SR and 48.2\% SPL), where motion trajectories provide temporal continuity and latent prediction offers geometric guidance.

\begin{table*}[t]
\centering
\caption{\textbf{Performance comparison across three navigation tasks.} 
For Image-Goal and Point-Goal, metrics are \textbf{SR} (\%) and \textbf{SPL} (\%). 
For No-Goal exploration, metrics are Explored \textbf{Area} ($\mathrm{m}^2$) and Episode \textbf{Time} (s). 
Best results are highlighted in \textbf{bold}.}
\label{tab:all_tasks_fixed_vertical}

\small
\renewcommand{\arraystretch}{0.95}

\begin{tabular*}{\textwidth}{@{\extracolsep{\fill}}l cc cc cc cc cc}
\toprule

\multirow{3}{*}{Method}
& \multicolumn{4}{c}{ClutterScenes}
& \multicolumn{4}{c}{InternScenes}
& \multicolumn{2}{c}{Average} \\

\cmidrule(lr){2-5}
\cmidrule(lr){6-9}
\cmidrule(lr){10-11}

& \multicolumn{2}{c}{Easy}
& \multicolumn{2}{c}{Hard}
& \multicolumn{2}{c}{Home}
& \multicolumn{2}{c}{Commercial}
& \multicolumn{2}{c}{} \\

\cmidrule(lr){2-3}
\cmidrule(lr){4-5}
\cmidrule(lr){6-7}
\cmidrule(lr){8-9}
\cmidrule(lr){10-11}

& $M_1$ & $M_2$
& $M_1$ & $M_2$
& $M_1$ & $M_2$
& $M_1$ & $M_2$
& $M_1$ & $M_2$ \\

\midrule

% ---------------- Task 1 ----------------
\multicolumn{11}{c}{\textbf{Task 1: Image-Goal Navigation} ($M_1$: SR $\uparrow$, $M_2$: SPL $\uparrow$)} \\
\cmidrule(lr){1-11}

GNM
& 26.9 & 25.9
& 20.0 & 19.2
& 7.7 & 7.3
& 10.5 & 10.3
& 16.3 & 15.7 \\

ViNT
& 15.0 & 13.9
& 16.5 & 15.2
& 6.2 & 6.0
& 12.9 & 12.2
& 12.6 & 11.8 \\

NoMaD
& 8.3 & 6.9
& 7.3 & 6.9
& 6.7 & 6.2
& 11.7 & 10.7
& 8.5 & 7.6 \\

NWM
& 12.7 & 12.0
& 8.7 & 8.4
& 4.0 & 3.9
& 6.9 & 6.6
& 8.1 & 7.7 \\

NavDP
& 49.7 & 48.4
& 47.5 & 46.3
& 29.8 & 27.7
& 46.6 & 43.3
& 43.4 & 43.4 \\

% \rowcolor[gray]{0.92}
Ours
& \textbf{66.1} & \textbf{64.5}
& \textbf{55.6} & \textbf{53.4}
& \textbf{30.7} & \textbf{29.4}
& \textbf{48.3} & \textbf{45.6}
& \textbf{50.2} & \textbf{48.2} \\

\midrule

% ---------------- Task 2 ----------------
\multicolumn{11}{c}{\textbf{Task 2: Point-Goal Navigation} ($M_1$: SR $\uparrow$, $M_2$: SPL $\uparrow$)} \\
\cmidrule(lr){1-11}

DD-PPO
& 0.0 & 0.0
& 0.0 & 0.0
& 0.4 & 0.4
& 5.3 & 5.2
& 1.4 & 1.4 \\

iPlanner
& 89.4 & 88.4
& 80.3 & 78.8
& 43.0 & 40.6
& 54.6 & 52.8
& 66.8 & 65.1 \\

ViPlanner
& 80.2 & 80.1
& 64.7 & 64.5
& 45.0 & 64.5
& 63.7 & 61.9
& 63.4 & 62.4 \\

NavDP
& 92.3 & 90.5
& 87.4 & 84.9
& 60.0 & 55.6
& 71.4 & 68.2
& 77.8 & 74.8 \\

% \rowcolor[gray]{0.92}
Ours
& \textbf{93.8} & \textbf{91.6}
& \textbf{88.9} & \textbf{86.4}
& \textbf{61.4} & \textbf{58.6}
& \textbf{77.4} & \textbf{75.3}
& \textbf{80.4} & \textbf{78.0} \\

\midrule

% ---------------- Task 3 ----------------
\multicolumn{11}{c}{\textbf{Task 3: No-Goal Exploration} ($M_1$: Area $\uparrow$, $M_2$: Time $\uparrow$)} \\
\cmidrule(lr){1-11}

GNM
& 42.8 & 17.1
& 29.6 & 12.5
& 19.8 & 30.8
& 19.2 & 28.1
& 27.8 & 22.1 \\

ViNT
& 55.3 & 20.3
& 46.6 & 18.9
& 20.8 & 31.5
& 22.4 & 27.2
& 36.2 & 24.4 \\

NoMaD
& 119.4 & 48.1
& 85.7 & 36.6
& 21.7 & 27.8
& 23.4 & 29.7
& 62.5 & 35.5 \\

NavDP
& 315.6 & 112.7
& 274.1 & 106.2
& \textbf{34.1} & \textbf{35.0}
& 43.9 & 36.2
& 167.2 & 72.5 \\

% \rowcolor[gray]{0.92}
Ours
& \textbf{320.4} & \textbf{115.2}
& \textbf{281.7} & \textbf{108.6}
& 27.9 & 30.3
& \textbf{45.5} & \textbf{40.2}
& \textbf{168.9} & \textbf{73.6} \\

\bottomrule
\end{tabular*}

\end{table*}

\begin{figure}[!t]
    \centering
    \includegraphics[width=\linewidth]{corl_2026_template_cameraready/figures/result.png}
    \caption{Qualitative comparison of Image-Goal navigation in IsaacSim. As illustrated by the 2D top-down trajectories and egocentric path projections, compared to NavDP, \texttt{WAM-Nav} shows stronger overall trajectory consistency and is smoother, and it can make better real-time obstacle avoidance predictions in advance. Notably, despite generating future states entirely in a compressed latent space, the decoded visual foresights maintain high fidelity to the ground truth (GT) observations.}
    \label{fig:result}
    \vspace{-5mm}
\end{figure}

\begin{table}[t]
\centering
\caption{\textbf{Ablation study of different components of \texttt{WAM-Nav} on Image-Goal navigation.}}
\label{tab:ablation_study_flattened_onedecimal_final}

\small
\renewcommand{\arraystretch}{1.1}

\begin{tabular*}{\columnwidth}{@{\extracolsep{\fill}}cc cc cc cc}
\toprule
\multicolumn{2}{c}{Method} &
\multicolumn{2}{c}{ClutterScenes (Avg.)} &
\multicolumn{2}{c}{InternScenes (Avg.)} &
\multicolumn{2}{c}{Overall Average} \\

\cmidrule(lr){1-2}
\cmidrule(lr){3-4}
\cmidrule(lr){5-6}
\cmidrule(lr){7-8}

Motion Traj. & Latent Pred. & SR & SPL & SR & SPL & SR & SPL \\
\midrule

$\times$ & $\times$ & 51.0 & 49.9 & 33.2 & 31.7 & 42.1 & 40.8 \\
$\times$ & $\checkmark$ & 57.4 & 55.5 & 34.0 & 32.2 & 45.7 & 43.8 \\
$\checkmark$ & $\times$ & 57.9 & 57.1 & 30.0 & 28.6 & 43.9 & 42.9 \\

% \rowcolor[gray]{0.92}
$\checkmark$ & $\checkmark$ &
\textbf{60.9} & \textbf{59.0} &
\textbf{39.5} & \textbf{37.5} &
\textbf{50.2} & \textbf{48.2} \\

\bottomrule
\end{tabular*}

% \vspace{-4mm}
\end{table}

\begin{figure}[!t]
    \centering
    \includegraphics[width=\linewidth]{corl_2026_template_cameraready/figures/realworld.png}
    \caption{Zero-shot deployment of \texttt{WAM-Nav} in real-world environments. (a) Visualization of \texttt{WAM-Nav} navigating across four representative indoor and outdoor scenes, with the first-person trajectory planning visualization shown in the bottom-right corner of each image. (b) Quantitative comparison of \texttt{WAM-Nav} and NavDP in each real-world scenario.}
    \label{fig:realworld}
    \vspace{-5mm}
\end{figure}

\subsection{Deployment on Real World}

% #模型效率

For \textbf{Q5}, we deploy \texttt{WAM-Nav} in a zero-shot manner on a real-world Unitree G1 humanoid robot equipped with an Intel RealSense D455 camera. Evaluation is conducted across four representative indoor and outdoor environments: a \textit{meeting room}, \textit{warehouse}, \textit{lobby}, and \textit{parking lot}, with 10 trials per scene. As shown in Fig.~\ref{fig:realworld}(a), \texttt{WAM-Nav} consistently predicts feasible trajectories across diverse spatial layouts and lighting conditions, enabling reliable obstacle avoidance and successful navigation in constrained environments. Compared with NavDP, \texttt{WAM-Nav} achieves a higher success rate and demonstrates stronger overall robustness in real-world deployment.

%==============================================================================

\section{Conclusion}
\label{sec:conclusion}

\textbf{Conclusion.}
This paper addresses the challenges of temporal trajectory inconsistency and myopic obstacle avoidance in embodied visual navigation by introducing \texttt{WAM-Nav}. Our key insight is to formulate future action prediction and perceptual foresight as a unified generative process. Built upon a shared DiT conditioned on a context-aware dual-stream prior, \texttt{WAM-Nav} jointly models robot motion dynamics and latent visual imagination, enabling the proactive generation of trajectories that are both dynamically consistent and geometrically safe. Extensive zero-shot experiments demonstrate that \texttt{WAM-Nav} achieves strong out-of-domain generalization across Image-Goal, Point-Goal, and No-Goal navigation in highly cluttered environments. Furthermore, a single trained policy generalizes effectively across diverse robot embodiments and successfully transfers to real-world humanoid platforms, highlighting its practical potential for embodied navigation.

\textbf{Limitations.}
During real-world deployment of \texttt{WAM-Nav}, we identify two primary failure modes. First, the limited camera height and field of view restrict perception of near-field obstacles, which may lead to inaccurate trajectory prediction or delayed avoidance. Second, the current policy does not explicitly model robot embodiment, causing mismatches between camera-level obstacle avoidance and full-body traversability. These limitations suggest two promising future directions: adaptive viewpoint control for dynamic perception adjustment, and embodiment-aware training with diverse robot morphologies to improve navigation robustness.
% 在真实环境中部署WAM-Nav时，我们发现两个主要问题会导致任务失败，一个是受相机高度及视角限制，由于看不到眼前障碍物，常会导致机器人路径轨迹规划错误或避障不及时，第二个是由于不考虑机器人自身形态，有时会出现相机成功避障但是机器人本体仍被障碍物阻止前进。因此，下一步可以研究一个舵机自适应算法，让相机像人眼一样能根据不通障碍物灵活调节观测视角，第二个可以在模型学习时，加入不同机器人本体数据，让导航能根据不同本体形态灵活避障。
%===============================================================================

\clearpage
% The acknowledgments are automatically included only in the final and preprint versions of the paper.
% \acknowledgments{If a paper is accepted, the final camera-ready version will (and probably should) include acknowledgments. All acknowledgments go at the end of the paper, including thanks to reviewers who gave useful comments, to colleagues who contributed to the ideas, and to funding agencies and corporate sponsors that provided financial support.}

%===============================================================================

% no \bibliographystyle is required, since the corl style is automatically used.
\bibliography{example}  % .bib

\newpage
\appendix
\section{Overview}

The supplementary material is organized as follows:
\begin{itemize}
    \item Section B provides additional details of the proposed method.
    \item Section C presents the model configurations and training details of \texttt{WAM-Nav}.
    \item Section D introduces the details of the baseline methods.
    \item Section E includes additional experimental results and ablation studies.
    \item Section F describes the hardware setup and deployment details for real-world experiments.
    \item Section G provides the visualization of a failure case.
\end{itemize}

\section{Method Details}
\label{appendix:method_details}

\subsection{Trajectory-Aware Motion History}
\label{appendix:ego_history_conversion}

While goal-modulated visual memory provides spatial context over $\mathcal{O}_t$, it does not explicitly encode kinematic momentum; the motion trajectory stream therefore maintains an action-accumulated ego-motion history $\mathcal{S}_t=\{s_{t-k+1}, \ldots, s_t\}$ (Sec.~\ref{sec:problem}). At deployment, each executed action $a_{t-1}$ updates the current planar pose $s_t$, which is appended to a length-$k$ sliding buffer and zero-padded at the front when the episode is shorter than $k$ steps. The buffer is then re-expressed in the current egocentric frame---re-centering on $s_t$, rotating translations by $\theta_t$, and wrapping heading differences---so that the stream receives a complete relative-motion sequence $\tilde{\mathcal{S}}_t=\{(\Delta x_i,\Delta y_i,\Delta\theta_i)\}_{i=t-k+1}^{t}$ rather than absolute scene coordinates. Algorithm~\ref{alg:ego_motion} details this online procedure.

\begin{algorithm}[h]
\caption{Online Ego-Motion History and Egocentric Conversion}
\begin{algorithmic}
\REQUIRE{window length $k$; executed action $a_{t-1}$; pose accumulator $s_{t-1}$; pose buffer $\mathcal{B}$}
\ENSURE{egocentric relative-motion sequence $\tilde{\mathcal{S}}_t=\{(\Delta x_i,\Delta y_i,\Delta\theta_i)\}_{i=t-k+1}^{t}$}
\STATE $s_t \leftarrow s_{t-1} + a_{t-1}$
\STATE Append $s_t$ to $\mathcal{B}$; if $|\mathcal{B}|>k$, remove the oldest pose
\IF{$|\mathcal{B}|<k$}
    \STATE $\mathcal{S}_t \leftarrow$ zero-pad the front of $\mathcal{B}$ to length $k$
\ELSE
    \STATE $\mathcal{S}_t \leftarrow \mathcal{B}$
\ENDIF
\STATE $s_t=(x_t,y_t,\theta_t) \leftarrow$ the last pose in $\mathcal{S}_t$
\STATE $\tilde{\mathcal{S}}_t \leftarrow \emptyset$
\FOR{each $s_i=(x_i,y_i,\theta_i)$ in $\mathcal{S}_t$}
    \STATE $\delta x \leftarrow x_i - x_t,\;\; \delta y \leftarrow y_i - y_t$
    \STATE $\Delta x \leftarrow \cos\theta_t \cdot \delta x + \sin\theta_t \cdot \delta y$
    \STATE $\Delta y \leftarrow -\sin\theta_t \cdot \delta x + \cos\theta_t \cdot \delta y$
    \STATE $\Delta\theta \leftarrow \mathrm{atan2}(\sin(\theta_i-\theta_t),\,\cos(\theta_i-\theta_t))$
    \STATE Append $(\Delta x,\Delta y,\Delta\theta)$ to $\tilde{\mathcal{S}}_t$
\ENDFOR
\end{algorithmic}
\label{alg:ego_motion}
\end{algorithm}

\subsection{Architecture of the Shared DiT Block}
\label{appendix:shared_dit}
As introduced in the \textbf{Coupled Action-Consequence Generation} subsection, the shared DiT predicts the joint flow-matching velocity fields $\hat{u}_A$ and $\hat{u}_Z$ from the interpolated action and latent visual states $(\mathbf{A}_{\tau}, \mathbf{Z}_{\tau})$, conditioned on the context prior $C$ and flow timestep $\tau$. Figure.~\ref{fig:shared_dit} illustrates a single DiT block, stacked $N$ times.

At each block, $\mathbf{A}_{\tau}$ and $\mathbf{Z}_{\tau}$ are tokenized into parallel action and image streams. The timestep embedding $\tau$ is injected into each stream via separate adaLN-Zero branches. The modulated tokens then pass through a shared self-attention layer, enabling early interaction between action and latent visual tokens. After the residual update, both streams perform cross-attention over the context prior $C$ using shared cross-attention and FFN weights with stream-specific modulation. The updated features are propagated through stacked blocks, and finally projected to $\hat{u}_A$ and $\hat{u}_Z$. This design enables consistent action--consequence generation through shared parameters and cross-modal coupling.

\begin{figure}[t]
    \centering
    \includegraphics[width=0.45\linewidth]{corl_2026_template_cameraready/figures/DiT.png}
    \caption{Architecture of the shared DiT block. Noised action and image tokens in latent-space are modulated through separate adaLN branches, coupled via self-attention, and grounded on $C$ through shared cross-attention and FFN layers.}
    \label{fig:shared_dit}
\end{figure}

\section{Implementation Details of WAM-Nav}
\label{appendix:implementation_details}

\begin{table}[t]
\centering
\caption{Key hyperparameters of \texttt{WAM-Nav}.}
\label{tab:hyperparameters}

\begin{tabular*}{\columnwidth}{@{\extracolsep{\fill}}lc|lc}
\toprule
\textbf{Parameter} & \textbf{Value} & \textbf{Parameter} & \textbf{Value} \\
\midrule
Hidden Dimension & 384 & Action Horizon ($H$) & 24 \\
Condition Length & 64 & Memory Window ($k$) & 8 \\
DiT Layers & 16 & DiT Attention Heads & 8 \\
Fusion Layers & 4 & Fusion Attention Heads & 8 \\
Causal Layers & 2 & Causal Attention Heads & 4 \\
Latent Grid Size & $16 \times 16$ & Latent Channels & 4 \\
Image Patches & 16 & Dropout & 0.1 \\
Batch Size & 256 & Learning Rate & $1.5 \times 10^{-4}$ \\
Weight Decay & $1\times10^{-4}$ & Warmup Ratio & 0.05 \\
$\lambda_{\mathrm{img}}$ & 0.25 & $\lambda_{\mathrm{align}}$ & 0.1 \\
Inference Steps & 10 & Visual Backbone & DINOv2 ViT-S/14 \\
\bottomrule
\end{tabular*}
\end{table}

We train \texttt{WAM-Nav} using a batch size of 512 and optimize the model with the AdamW optimizer. The peak learning rate is set to $1 \times 10^{-4}$ with a weight decay of $1 \times 10^{-4}$, following a cosine annealing schedule with a 5\% linear warmup. During training, the visual backbone (DINOv2 ViT-S/14) and the pre-trained Stable Diffusion VAE are kept frozen to preserve their general-purpose representations, while the goal image encoder, fusion decoder, causal motion encoder, and the shared DiT are trained from scratch. The loss weighting coefficients in Eq.~\ref{eq:l_total} are empirically set to $\lambda_{\mathrm{img}}=0.05$ and $\lambda_{\mathrm{align}}=0.01$. At deployment, the flow-matching ODE solver runs for $N_{\mathrm{steps}}=10$ Euler integration steps, which balances high-quality trajectory generation with real-time responsiveness. The detailed hyperparameter configurations are summarized in Table~\ref{tab:hyperparameters}.

\section{Baseline Details}
\label{appendix:baseline_details}
To comprehensively evaluate \texttt{WAM-Nav}, we compare against representative baselines spanning reinforcement learning, reactive planning, foundation models, and generative world models.

\begin{itemize}
    \item \textbf{DD-PPO}~\citep{wijmans2020ddppo}: A large-scale distributed reinforcement learning framework for Point-Goal navigation, trained over billions of frames in Habitat and widely adopted as a strong mapless navigation baseline.

    \item \textbf{iPlanner}~\citep{yang2023iplanner}: A high-frequency local planner that directly predicts multi-step continuous trajectories through spline regression for reactive obstacle avoidance.

    \item \textbf{ViPlanner}~\citep{roth2024viplanner}: An end-to-end semantic planner that predicts 2D collision cost maps from monocular observations for safe trajectory generation.

    \item \textbf{GNM}~\citep{shah2023gnm}: A general navigation foundation model trained on large-scale cross-embodiment datasets for zero-shot transfer and low-level action prediction.

    \item \textbf{ViNT}~\citep{shah2023vint}: A Transformer-based navigation policy that extends GNM by jointly modeling heterogeneous visual navigation data with improved sequence representation.

    \item \textbf{NoMaD}~\citep{sridhar2024nomad}: A diffusion-based action policy that introduces goal masking to unify goal-conditioned navigation and goal-free exploration.

    \item \textbf{NWM}~\citep{bar2025nwm}: A generative navigation world model based on conditional DiT that predicts future egocentric observations for foresight-driven planning.

    \item \textbf{NavDP}~\citep{cai2025navdp}: A diffusion-based navigation policy that leverages privileged simulation priors for zero-shot Sim-to-Real transfer.
\end{itemize}

\section{Additional Experimental Results}
\label{appendix:additional_results}
\subsection{Ablation Studies}

To systematically investigate the contribution of each key component in \texttt{WAM-Nav}, we conduct a comprehensive ablation study on the Image-Goal navigation task. Specifically, we analyze three aspects: the historical memory window size ($k$), the dual-stream information fusion design in DSCC, and the coupling strategy between action policy generation and latent visual imagination.

For the memory mechanism, we evaluate window sizes of $k \in \{4, 8, 16\}$. For DSCC, we compare our goal-modulated memory gating design against a baseline that directly concatenates visual, kinematic, and goal tokens without explicit goal injection, denoted as \textit{w/o GI}. For action-imagination coupling, we further study two architectural variants: (1) a \textit{Decoupled DiT}, where the policy and imagination branches are parameterized independently without weight sharing; and (2) a \textit{Partially Shared DiT}, where only the self-attention layers are shared, while the cross-attention and feed-forward network (FFN) layers remain separate. Quantitative results are summarized in Table~\ref{tab:ablation_combined}.

As shown in Table~\ref{tab:ablation_combined}, both insufficient and excessive temporal context negatively affect navigation performance. Figure.~\ref{fig:ablation_traj} further visualizes how different memory window sizes influence trajectory execution. With a shorter window ($k=4$), the agent is able to complete simple navigation tasks, but the generated trajectories are less smooth and exhibit noticeable instability. This suggests that limited temporal context constrains the model to short-horizon planning, weakening its ability to perform long-range geometric reasoning. In contrast, when using a larger memory window ($k=16$), the model becomes overly dependent on historical motion patterns, which reduces its responsiveness to sudden geometric changes in the environment and leads to more frequent collisions.

Removing explicit goal injection (\textit{w/o GI}) also causes a clear performance drop, highlighting the importance of goal-aware modulation in constructing effective contextual priors. Without explicit goal conditioning, the model exhibits weaker directional guidance and reduced task success.

Finally, the coupling mechanism between policy generation and latent imagination plays a critical role in overall performance. Although the \textit{Decoupled DiT} variant provides basic predictive visual foresight and already surpasses the reactive NavDP baseline, the absence of parameter sharing limits information exchange between trajectory planning and future visual prediction, resulting in weaker mutual regularization and structural inconsistency. The \textit{Partially Shared DiT} alleviates this issue to some extent, but still fails to fully align action generation with predicted future observations. In contrast, our fully shared DiT architecture tightly couples the two branches throughout all transformer layers, enabling latent visual imagination to directly regularize policy generation under unified spatiotemporal representations. This strong coupling imposes more consistent geometric constraints on the action space, suppresses dynamically infeasible trajectories, and ultimately leads to the best navigation performance and trajectory stability.

\begin{table*}[t]
\centering
\small
\renewcommand{\arraystretch}{1.15}
\caption{Ablation studies on memory window size, DSCC module, and Policy \& Imagine architecture.}
\label{tab:ablation_combined}

\begin{tabular*}{\textwidth}{@{\extracolsep{\fill}}lcc|lcc|lcc}
\toprule
\multicolumn{3}{c|}{\textbf{Memory Window Size}} &
\multicolumn{3}{c|}{\textbf{DSCC}} &
\multicolumn{3}{c}{\textbf{Policy \& Imagine}} \\
\cmidrule(r){1-3}
\cmidrule(r){4-6}
\cmidrule(l){7-9}

\textbf{Setting} & SR ($\uparrow$) & SPL ($\uparrow$) &
\textbf{Variant} & SR ($\uparrow$) & SPL ($\uparrow$) &
\textbf{Variant} & SR ($\uparrow$) & SPL ($\uparrow$) \\
\midrule

$k=4$ & 46.6 & 44.1 &
-- & -- & -- &
Decoupled DiT & 45.9 & 44.6 \\

$k=8$ (Ours) & \textbf{50.2} & \textbf{48.2} &
w/o GI & 44.5 & 42.2 &
Partially Shared DiT & 47.5 & 45.4 \\

$k=16$ & 44.9 & 42.3 &
Ours & \textbf{50.2} & \textbf{48.2} &
Ours & \textbf{50.2} & \textbf{48.2} \\

\bottomrule
\end{tabular*}
\end{table*}

\begin{figure}[t]
    \centering
    \includegraphics[width=0.75\linewidth]{corl_2026_template_cameraready/figures/ablation_result.png}
    \caption{Visualization of the effect of different memory window sizes on navigation trajectory planning.}
    \label{fig:ablation_traj}
\end{figure}
% 1、复杂任务应对能力
% 2、跨本体适应能力
% \textbf{Q4:} Can a single \texttt{WAM-Nav} policy generalize across diverse robot embodiments without retraining? 

% For \textbf{Q4}, Table~\ref{tab:embodiment_fixed} evaluates the zero-shot adaptability of a single \texttt{WAM-Nav} policy across diverse robot embodiments, including the wheeled Dingo and Unitree humanoid robots (G1 and H2) in IsaacSim. The results demonstrate strong cross-embodiment generalization. On G1, \texttt{WAM-Nav} consistently outperforms NavDP across all tasks, achieving 42.9\% SR on Image-Goal navigation. On H2, although NavDP performs slightly better on Point-Goal navigation (52.7\% vs. 47.9\% SR), \texttt{WAM-Nav} achieves higher success rates in Image-Goal and broader exploration coverage in No-Goal scenarios. These results suggest that \texttt{WAM-Nav} effectively captures embodiment-agnostic navigation priors, enabling robust transfer across heterogeneous platforms without retraining.
\subsection{Comparison with NavDP}
\label{subsec:compared_navdp}

We first evaluate the zero-shot adaptability of a single \texttt{WAM-Nav} policy across diverse robot embodiments without any parameter retraining. As summarized in Table~\ref{tab:embodiment_fixed}, experiments are conducted on a wheeled Dingo platform together with Unitree G1 and H2 humanoid robots in the IsaacSim environment. The quantitative results demonstrate strong cross-embodiment generalization. In particular, on the G1 platform, \texttt{WAM-Nav} consistently outperforms NavDP across all evaluated task settings, achieving a $42.9\%$ SR on Image-Goal navigation. On the more challenging H2 humanoid platform, \texttt{WAM-Nav} also achieves higher performance in all navigation tasks. These results suggest that \texttt{WAM-Nav} effectively learns embodiment-agnostic navigation priors, enabling stable policy transfer across robots with substantially different kinematic structures without retraining.

To further characterize the capabilities of our framework, we compare \texttt{WAM-Nav} and NavDP across navigation tasks of varying difficulty under the Image-Goal setting. As shown in Table~\ref{tab:difficulty}, navigation episodes are divided into Easy, Medium, and Hard subsets according to the ground-truth path length. Across all difficulty levels, \texttt{WAM-Nav} consistently outperforms NavDP, with the performance gap increasing as task difficulty grows. Notably, in the Hard subset, \texttt{WAM-Nav} reaches $25.0\%$ SR and $23.8\%$ SPL, surpassing NavDP by $7.6\%$ and $2.5\%$, respectively.

This growing advantage under long-horizon navigation highlights the benefit of our coupled action-consequence generation framework. Over extended trajectories, reactive diffusion policies such as NavDP tend to accumulate execution drift and are more vulnerable to local geometric traps due to the lack of predictive perceptual feedback. In contrast, \texttt{WAM-Nav} jointly models action generation and latent visual imagination, allowing predicted future observations to continuously regularize trajectory planning. This coupled predictive mechanism enforces long-range spatiotemporal consistency, reduces myopic behavioral drift, and enables the policy to maintain stable and goal-directed navigation even in long-distance and cluttered environments.

\begin{table}[t]
\centering
\caption{\textbf{Generalization performance across different embodiments.}}
\label{tab:embodiment_fixed}

\small
\renewcommand{\arraystretch}{1.05}

\begin{tabular*}{\columnwidth}{@{\extracolsep{\fill}}cl cc cc cc}
\toprule

\multirow{2}{*}{Embodiment}
& \multirow{2}{*}{Method}
& \multicolumn{2}{c}{Image-Goal}
& \multicolumn{2}{c}{Point-Goal}
& \multicolumn{2}{c}{No-Goal} \\

\cmidrule(lr){3-4}
\cmidrule(lr){5-6}
\cmidrule(lr){7-8}

&
& SR ($\uparrow$) & SPL ($\uparrow$)
& SR ($\uparrow$) & SPL ($\uparrow$)
& Area ($\uparrow$) & Time ($\uparrow$) \\

\midrule

% Dingo
\multirow{2}{*}{Dingo}
& NavDP
& 43.4 & 43.4
& 77.8 & 74.8
& 167.2 & 72.5 \\

& Ours
& \textbf{50.2} & \textbf{48.2}
& \textbf{80.4} & \textbf{78.0}
& \textbf{168.9} & \textbf{73.6} \\

\midrule

% G1
\multirow{2}{*}{G1}
& NavDP
& 37.2 & 35.6
& 60.3 & 56.8
& 133.0 & 65.2 \\

& Ours
& \textbf{42.9} & \textbf{42.5}
& \textbf{64.8} & \textbf{63.1}
& \textbf{147.6} & \textbf{68.3} \\

\midrule

% H2
\multirow{2}{*}{H2}
& NavDP
& 32.7 & 31.3
& 47.9 & 46.5
& 142.4 & 67.6 \\

& Ours
& \textbf{35.5} & \textbf{34.9}
& \textbf{52.7} & \textbf{50.5}
& \textbf{145.4} & \textbf{70.6} \\

\bottomrule
\end{tabular*}

\end{table}

\begin{table}[t]
\centering
\caption{Performance across navigation difficulty levels.}
\label{tab:difficulty}
\small
\setlength{\tabcolsep}{7pt}
\renewcommand{\arraystretch}{1.05}

\begin{tabular*}{\columnwidth}{@{\extracolsep{\fill}}l cc cc cc cc}
\toprule
\multirow{2}{*}{Method}
& \multicolumn{2}{c}{Easy}
& \multicolumn{2}{c}{Medium}
& \multicolumn{2}{c}{Hard}
& \multicolumn{2}{c}{Average} \\
\cmidrule(lr){2-3}
\cmidrule(lr){4-5}
\cmidrule(lr){6-7}
\cmidrule(lr){8-9}
& SR ($\uparrow$) & SPL ($\uparrow$)
& SR ($\uparrow$) & SPL ($\uparrow$)
& SR ($\uparrow$) & SPL ($\uparrow$)
& SR ($\uparrow$) & SPL ($\uparrow$) \\
\midrule
NavDP & 46.2 & 45.7 & 33.8 & 35.6 & 17.4 & 21.3 & 43.4 & 43.4 \\
Ours  & \textbf{52.5} & \textbf{50.5}
      & \textbf{42.4} & \textbf{40.5}
      & \textbf{25.0} & \textbf{23.8}
      & \textbf{50.2} & \textbf{48.2} \\
\bottomrule
\end{tabular*}
\end{table}

\begin{figure}[b]
    \centering
    \includegraphics[width=0.75\linewidth]{corl_2026_template_cameraready/figures/G1.png}
    \caption{Robot Setup. (a) Unitree G1. (b) Detailed view of the hardware components.}
    \label{fig:G1}
\end{figure}

\section{Real-World Deployment Details}
\label{appendix:realworld_details}
\subsection{Hardware setup}
As shown in Figure.~\ref{fig:G1}, our experimental platform is built on the Unitree G1 humanoid robot. The robot is equipped with an Intel RealSense D455 depth camera mounted on a pan-tilt servo to capture RGB-D observations from fixed viewpoints. A backpack-mounted onboard computer equipped with an RTX 4060 GPU performs real-time perception and control. A power transformer provides stable power supply for the onboard devices. This hardware setup enables real-time perception, computation, and execution for our experiments.

\subsection{Deployment Details}
During deployment, the Intel RealSense D455 depth camera is mounted on the pan-tilt servo and configured with a 20° downward pitch angle to improve the observation of ground-level obstacles. The camera captures RGB-D observations at 30 Hz.
The WAM-Nav model is deployed as a server on the backpack-mounted onboard computer equipped with an RTX 4060 GPU for online inference. The model occupies approximately 1.3 GB of GPU memory, and the inference frequency is set to 1 Hz.
The client module is also deployed on the onboard computer. It is responsible for forwarding the incoming observations to the server and receiving the predicted outputs. The predicted trajectories are tracked using Model Predictive Control (MPC), which further converts them into velocity commands for the robot motors. During experiments, the control loop runs at 50 Hz, with each command applied for 0.1 s.

\section{Additional Failure Case Analysis}
\label{appendix:failure_analysis}
\begin{figure}[!h]
    \centering
    \includegraphics[width=0.75\linewidth]{corl_2026_template_cameraready/figures/failure.png}
    \caption{Failure Cases. The bottom-left subfigure shows the first-person view of the planned trajectory}
    \label{fig:failure}
\end{figure}
In Figure~\ref{fig:failure}, we visualize several typical failure cases. Case (a) shows that when the robot is close to a low obstacle, the limited field of view results in a weak perception of the obstacle, causing the robot to overlook it and eventually collide. Another common type of failure is illustrated in (b) and (c): when obstacles are located on the left or right side of the robot, the robot’s body shape is not taken into account. As a result, trajectory planning only ensures that the camera can pass through, leading to collisions between the robot and the obstacles.
\end{document}
